\chapter{Formación y Desarrollo}

\section{Programa de formación utilizado (identificación, diseño e implantación)}

La gestión del aprendizaje en la organización se desarrolla siguiendo un ciclo técnico que busca transformar las capacidades de los empleados para alinearlas con la competitividad del sector bancario. Este proceso se divide en las fases que se detallan a continuación.  

\section{Programa de formación utilizado (identificación, diseño e implantación)}

En CaixaBank, la mejora continua de su plantilla se estructura partiendo de una distinción teórica fundamental que la entidad aplica en su día a día: la diferencia entre formación y desarrollo. Por un lado, la entidad utiliza la formación como un proceso reactivo centrado en el trabajo actual; por ejemplo, instruyendo a los empleados provenientes de la fusión con Bankia en el uso del software propio de la entidad o actualizando conocimientos normativos inmediatos. Por otro lado, aplica el desarrollo como un proceso proactivo y a largo plazo, preparando a empleados talentosos mediante planes de carrera para que asuman futuras posiciones de liderazgo o perfiles altamente tecnificados.

Para llevar a cabo ambos procesos, el Área de Personas y Organización de CaixaBank sigue un ciclo técnico estructurado en las siguientes fases:

\subsection{Identificación de necesidades de formación}

Esta etapa inicial busca diagnosticar las carencias del personal para justificar la inversión en aprendizaje. Siguiendo los principios teóricos, los responsables de recursos humanos de CaixaBank deben dar respuesta a tres preguntas fundamentales antes de planificar cualquier acción:

\begin{itemize}
    \item ¿Qué conocimientos y habilidades concretas en el desempeño del trabajo mejorarían la productividad laboral?
    \item ¿Cuál es el público objetivo?
    \item ¿Puedo desarrollar objetivos mensurables para esos conocimientos y habilidades?
\end{itemize}

Para responder a estas cuestiones con precisión, la entidad evalúa sus necesidades estratificando la información en tres niveles:

\begin{itemize}
    \item Nivel de organización: Se evalúan las necesidades derivadas de la estrategia global del banco. Por ejemplo, observando el Plan Estratégico 2025-2027, la dirección ha identificado que toda la organización necesita formarse urgentemente en inteligencia artificial generativa, finanzas sostenibles y atención a los nuevos ecosistemas vinculados al envejecimiento poblacional (segmento sénior).
    \item Nivel de tareas: Apoyándose en el departamento de Secretaría Técnica de Negocios y consultoras como Deloitte, se analiza qué exige cada puesto concreto. Así, detectan que el puesto de Gestor de Banca Privada requiere obligatoriamente habilidades analíticas complejas y certificaciones financieras reguladas.
    \item Nivel de individuo: Utilizando su Sistema de Información de Recursos Humanos (la plataforma en la nube SAP SuccessFactors), cruzan los datos de la evaluación del desempeño de cada trabajador con los requisitos de su puesto. Esto permite identificar de forma automatizada las brechas entre el rendimiento actual del empleado y su potencial, generando alarmas sobre qué formación específica necesita cada individuo.
\end{itemize}

\subsection{Diseño de los programas de formación}

Una vez que SAP SuccessFactors y los directores de área han detectado las necesidades, se procede al diseño del programa. En esta etapa, la teoría exige concretar los siguientes parámetros:

\begin{itemize}
    \item ¿Qué tipo de formación impartir?
    \item ¿Qué grado de aprendizaje alcanzar?
    \item ¿Para quién?
    \item ¿Quién la va a impartir?
\end{itemize}

CaixaBank diseña los contenidos estructurándolos en diferentes grados de aprendizaje. Para niveles iniciales, diseñan programas de acogida (onboarding) centrados en el saber; para niveles medios, diseñan formaciones en habilidades comerciales y trabajo en equipo orientadas al poder hacer; y para el desarrollo directivo, se enfocan en habilidades conceptuales.

Respecto a quién imparte la formación, la entidad utiliza un enfoque mixto. Como formadores internos, CaixaBank emplea a sus propios supervisores y a empleados senior mediante programas de mentoría, algo esencial para transmitir la cultura organizativa y orientar a los nuevos. Como formadores externos, cuando el grado de aprendizaje requerido es muy elevado y técnico, la entidad externaliza el servicio. Por ejemplo, se apoya en alianzas con la UPF Barcelona School of Management y otras instituciones académicas para impartir los programas que otorgan a sus empleados las certificaciones europeas MiFID II o EFPA, obligatorias para el asesoramiento financiero.

\subsection{Implantación de los programas de formación}

La última fase operativa requiere decidir el entorno y los canales tecnológicos que se utilizarán, dando respuesta a dos cuestiones clave:

\begin{itemize}
    \item ¿Dónde se realizará la formación?
    \item ¿Qué medios se utilizarán?
\end{itemize}

Dada la inmensa capilaridad de su red de oficinas (miles de sucursales en todo el territorio), CaixaBank combina metodologías para asegurar que la formación llegue a toda la plantilla de manera eficiente:

\begin{itemize}
    \item Métodos en el puesto de trabajo: Se aplican mediante la rotación de puestos (moving interno) y el aprendizaje directo. Por ejemplo, la asignación de un compañero experimentado (buddy) a los nuevos gestores comerciales permite que el trabajador adquiera la experiencia resolviendo incidencias reales con clientes, lo que asegura una transferencia directa del conocimiento sin abandonar la sucursal.
    \item Métodos fuera del lugar de trabajo (y entornos virtuales): La principal herramienta de implantación masiva de CaixaBank es su plataforma corporativa de aprendizaje denominada Virtaula. Se trata de un Learning Management System (LMS) donde la entidad aloja todo su e-learning. Dentro de Virtaula, el conocimiento está segmentado en escuelas especializadas (como la Escuela de Riesgos o la Escuela de Liderazgo). 
\end{itemize}

Además, en la implantación actual, CaixaBank está utilizando licencias de Microsoft Copilot integradas en sus sistemas. Esta inteligencia artificial actúa en la fase de implantación recomendando a cada empleado, de forma personalizada, qué módulos o píldoras formativas de Virtaula debe cursar a continuación según su progreso, permitiendo un aprendizaje a ritmo individual que reduce drásticamente los costes logísticos frente a los antiguos cursos presenciales.  

\section{Evaluación de los programas formativos}

La fase de evaluación constituye el cierre del ciclo de formación en CaixaBank y tiene como propósito fundamental determinar si las actividades realizadas han satisfecho los objetivos establecidos inicialmente. Este proceso no se limita a una comprobación administrativa del cumplimiento de los cursos, sino que busca medir si la inversión ha servido para mejorar el rendimiento presente y futuro de los empleados de la entidad.

Para que este análisis sea riguroso, el departamento de recursos humanos de la organización aplica un modelo que busca dar respuesta a los siguientes niveles de control definidos en la teoría:

\begin{itemize}
\item Reacción: ¿indican las reacciones iniciales de las personas ante la experiencia formativa que el aprendizaje es relevante e inmediatamente aplicable?
\item Aprendizaje: ¿han adquirido los conocimientos y habilidades planificadas?
\item Conducta: ¿se aplica lo aprendido en el trabajo?
\item Rendimiento: ¿qué resultados se aprecian en el negocio?
\item Impacto: ¿cuál es el coste-beneficio para la empresa?
\end{itemize}

Para llevar a la práctica estos niveles de evaluación en una red tan extensa, CaixaBank utiliza diversas herramientas tecnológicas y metodológicas integradas en su ecosistema digital:

Para medir la reacción y el aprendizaje, la entidad aprovecha las capacidades analíticas de su plataforma Virtaula. Al finalizar cada módulo formativo, especialmente en los itinerarios de e-learning, el sistema requiere que el empleado complete cuestionarios de satisfacción donde se evalúa la utilidad percibida de los contenidos. Asimismo, el aprendizaje se valida mediante test y pruebas de conocimientos integradas en el propio entorno virtual. Estas pruebas son obligatorias en ámbitos de cumplimiento normativo, como la formación en blanqueo de capitales, o para obtener las certificaciones financieras externas necesarias para operar en banca privada.

La evaluación de la conducta se realiza de manera más cualitativa y a medio plazo a través de la gestión del rendimiento. La organización lleva a cabo seguimientos periódicos, usualmente a los seis meses y al año de una formación relevante o de un cambio de puesto. En estas etapas, se analiza si el trabajador aplica efectivamente lo aprendido para agilizar procesos o mejorar la atención al cliente. En este nivel, la entrevista de evaluación con el responsable directo es la herramienta clave para confirmar si las nuevas habilidades técnicas o actitudes se han incorporado con éxito al día a día de la sucursal.

Finalmente, el rendimiento y el impacto se miden contrastando los resultados de la formación con las metas estratégicas del banco. Utilizando People Analytics dentro de su sistema SAP, CaixaBank monitoriza si los programas de reskilling y reciclaje profesional han contribuido a alcanzar indicadores de negocio como el aumento del margen de intereses o la mejora de la eficiencia operativa proyectada en el Plan Estratégico 2025-2027. Este análisis final permite a la dirección determinar si el coste de los programas ha generado un beneficio tangible que justifique el presupuesto invertido en el capital humano.

\section{Análisis de tiempo y coste de los programas}

En CaixaBank, la gestión del tiempo y el coste de la formación no se analiza de forma aislada, sino como una pieza fundamental del proceso de dimensionamiento que determina la estructura global de costes del banco. La entidad entiende que cualquier actividad formativa conlleva una inversión de recursos que debe ser optimizada para asegurar que el aprendizaje sea eficiente y no interfiera excesivamente en la operatividad diaria de las oficinas.

\subsection{Gestión y optimización del tiempo formativo}

La variable tiempo es crítica en el sector bancario, donde la carga de trabajo operativa es elevada. Según la teoría, los métodos de formación en el puesto de trabajo, aunque efectivos, suelen requerir mucho tiempo y pueden interferir con el rendimiento real. CaixaBank gestiona este riesgo mediante las siguientes estrategias:

\begin{itemize}
    \item Uso de Virtaula para el aprendizaje asíncrono: Al emplear un sistema de formación por ordenador y vídeo interactivo, la entidad permite el aprendizaje a ritmo individual. Esto reduce el tiempo que los empleados deben pasar fuera de sus puestos en cursos presenciales y elimina los tiempos de desplazamiento.
    \item Microaprendizaje o Instant Learning: La plataforma permite consumir soluciones instantáneas a carencias profesionales mediante píldoras informativas breves. Esto responde a la necesidad de formación inmediata sin comprometer largas jornadas laborales.
    \item Planificación estratégica del tiempo: La entidad utiliza métricas de seguimiento para monitorizar el tiempo de contratación y el tiempo de integración de nuevos empleados. El objetivo es que la curva de aprendizaje sea lo más corta posible para que el profesional alcance su nivel de productividad óptimo en el menor plazo de tiempo.
\end{itemize}

\subsection{Análisis y estructura de costes de la formación}

Aunque la organización no publica una cifra cerrada dedicada exclusivamente a la formación, su estrategia de costes se integra en la inversión masiva de más de 5.000 millones de euros destinada a tecnología y transformación para el periodo 2025-2027. La entidad diferencia entre diversos tipos de costes según la metodología empleada:

\begin{itemize}
    \item Costes de formación interna: Son los derivados del uso de supervisores y compañeros como formadores. Aunque estos métodos son teóricamente más baratos por no requerir instalaciones externas, CaixaBank asume el coste indirecto del tiempo que estos profesionales dejan de dedicar a sus tareas comerciales para actuar como mentores.
    \item Costes de formación externa y especializada: La entidad presupuesta partidas específicas para colaborar con asesores externos y universidades. Estos programas, como los necesarios para las certificaciones financieras, tienen un coste directo elevado pero garantizan un nivel de aprendizaje superior y el cumplimiento de la normativa legal.
    \item Inversión en infraestructuras digitales: El desarrollo y mantenimiento de sistemas como SAP SuccessFactors y Virtaula supone un coste de preparación muy alto. Sin embargo, una vez implantados, permiten formar a grandes grupos de personas (los más de 45.000 empleados del grupo) de forma mucho más barata y accesible que los cursos reglados tradicionales.
\end{itemize}

Finalmente, el banco utiliza herramientas de Learning Analytics para evaluar el impacto económico. Esto permite identificar patrones de aprendizaje y adaptar los itinerarios, asegurando que cada euro invertido en formación genere un retorno medible en términos de mejora de capacidades y cumplimiento de los objetivos estratégicos del grupo.

\section{Tendencias en formación: E-learning, Método 70-20-10 y Universidades Corporativas}

CaixaBank ha transformado su modelo de aprendizaje adoptando tendencias que permiten una formación más flexible, colaborativa y alineada con los objetivos de desarrollo a largo plazo. La entidad no solo busca corregir carencias inmediatas a través de cursos tradicionales, sino que utiliza metodologías avanzadas para potenciar el talento de forma continua.

\subsection{E-learning y entornos virtuales de aprendizaje}

La formación en línea es el eje central de la estrategia de aprendizaje de la entidad. Basada en el concepto de aprendizaje 3.0, la organización fomenta un entorno autónomo y colaborativo que elimina las barreras geográficas y temporales. La herramienta principal para esta tendencia es Virtaula, el Learning Management System del banco. Esta plataforma permite la gestión de contenidos digitales y facilita el microaprendizaje o instant learning, proporcionando soluciones rápidas a necesidades profesionales puntuales.

Además, la entidad está incorporando analítica de aprendizaje (learning analytics) y big data para identificar patrones de comportamiento en los estilos de estudio de sus empleados. Esto permite crear itinerarios de aprendizaje hiper-personalizados en la nube. La reciente integración de inteligencia artificial generativa, como Microsoft Copilot, refuerza esta tendencia al asistir a los empleados en la búsqueda de contenidos relevantes dentro de la plataforma y sugerir formaciones de manera proactiva según el perfil técnico o comercial del usuario.

\subsection{El modelo de aprendizaje 70-20-10}

CaixaBank aplica de manera práctica la metodología 70-20-10, la cual establece que el aprendizaje efectivo se produce a través de tres vías diferenciadas:

En primer lugar, el 70 por ciento del aprendizaje es experiencial. Se obtiene mediante la práctica diaria, la resolución de problemas en la oficina y la participación en nuevos proyectos o tareas. La entidad potencia este bloque a través de la rotación de puestos, permitiendo que el trabajador asuma retos en diferentes áreas para adquirir una visión global del negocio.

En segundo lugar, el 20 por ciento corresponde a la interacción social. Para cubrir este porcentaje, la organización utiliza sistemas de coaching y mentoring. Un ejemplo claro es la asignación de un compañero o mentor (buddy) a las nuevas incorporaciones, facilitando que el conocimiento se transmita de forma orgánica a través del debate y el trabajo en equipo dentro de la propia sucursal.

En tercer lugar, el 10 por ciento restante es el aprendizaje formal. Este se canaliza principalmente a través de Virtaula mediante cursos, talleres y seminarios técnicos. Aunque es el porcentaje menor, es fundamental para asentar las bases teóricas y cumplir con las certificaciones normativas obligatorias del sector bancario.

\subsection{Estructura de Universidad Corporativa}

Aunque la entidad no utiliza siempre el término comercial de universidad corporativa, su estructura interna de escuelas especializadas funciona bajo este concepto teórico. Se trata de estructuras diseñadas para asegurar que el aprendizaje esté directamente conectado con la estrategia y los procesos del banco.

CaixaBank cuenta con diversas escuelas de conocimiento, como la Escuela de Riesgos, la Escuela de Liderazgo o la Escuela Comercial. Estas unidades son las encargadas de transmitir no solo los conocimientos técnicos necesarios para el sector, sino también los valores, la cultura y la historia de la organización. Este modelo asegura que la formación sea coherente en todo el grupo y que los programas de desarrollo preparen a los mandos medios y superiores para los desafíos estratégicos del futuro.

\subsection{Gamificación y nuevas dinámicas}

La organización está implementando estrategias de juego en entornos profesionales para aumentar el compromiso y reforzar habilidades como la creatividad o la comunicación interna. A través de dinámicas interactivas y test en formato de retos, la entidad busca que la formación continua tenga un impacto positivo en la motivación de los trabajadores. Estas técnicas se utilizan tanto en los procesos de integración de nuevos empleados como en los programas de actualización de competencias para perfiles directivos, facilitando un aprendizaje más dinámico y efectivo.